% SHARD-ID: AQM-94-TRANSLATION-ALGEBRA-SPEC
% SHARD-TITLE: Translation Haah Data versus Bare Spectrum
% SHARD-SUMMARY: Isolates the most restrictive Haah layer: translation group algebras, antipodes, identity-coefficient traces, and finite periodic quotients.
% SHARD-SUMMARY: Explains why the translation basis of \(\mathbb F_p[\Lambda]\) is not the Zariski point set of \(\Spec\mathbb F_p[\Lambda]\).
% SHARD-SUMMARY: Gives examples and a checklist for recognizing when a scheme proposal has reached the translation Haah layer.
% SHARD-KEYWORDS: Haah, translation algebra, Spec, Laurent module, group algebra, antipode, toric code, finite periodic quotient

\section{Translation Haah Data versus Bare Spectrum}
\label{sec:translation-haah-data-versus-bare-spectrum}

\subsection{Status and Source Anchors}

This shard is a local synthesis, not a new classification theorem.  Its Haah
anchors are \path{references/lattice_stabilizer_codes/SOURCES.md} and the
local TeX locators
\path{references/lattice_stabilizer_codes/haah_1204_1063_source/alg-code-hamiltonian-v6.tex}
lines 169--183, 416--432, 514--544, and 634--667;
\path{references/lattice_stabilizer_codes/haah_1305_6973_source/alg-theory.tex}
lines 700--716 and 1515--1567; and
\path{references/lattice_stabilizer_codes/haah_1812_11193_source/cph2d-arXiv3.tex}
lines 321--365, 409--453, and 505--511.  The torus comparison is AQM-90,
AQM-91, and \texttt{CONVENTIONS.md} \texttt{(az)}--\texttt{(ba)}.  The
residue-spectrum and product-field warnings are AQM-40, AQM-81, AQM-36,
AQM-37, and conventions \texttt{(ad)}, \texttt{(aa)}, and \texttt{(ab)}.

\subsection{Lamport Dependency Skeleton}
\[
\begin{gathered}
D1:(G,P,\lambda,\sigma)\text{ is the permissive pre-Haah datum over }A.\\
D2:\Spec A\text{ adds sheaf/support/fibre geometry but not sites.}\\
D3:\text{finite Weyl data adds a self-dual label convention.}\\
D4:R=\mathbb F_p[\Lambda],\ \bar{(\cdot)},\ \operatorname{tr},\ \dagger
   \text{ are the translation layer.}\\
D5:P=R^{2q},\quad
\lambda_q=\begin{pmatrix}0&I_q\\-I_q&0\end{pmatrix},\quad
\epsilon=\sigma^\dagger\lambda_q.\\
L1:\sigma^\dagger\lambda_q\sigma=0,\qquad
L2:\ker\epsilon=\operatorname{im}\sigma\text{ only when sourced.}
\end{gathered}
\]

\subsection{From Permissive to Restrictive}

Following \texttt{CONVENTIONS.md} \texttt{(ba)}, the permissive layer is a
\emph{pre-Haah datum} \((G,P,\lambda,\sigma)\) over a commutative ring \(A\):
modules, a chosen alternating or symplectic form when available, and a
generator map.  Here \(\epsilon=\sigma^\dagger\lambda\) and
\(\ker\epsilon/\operatorname{im}\sigma\) are only local definitions of a
logical-defect module, not a physical logical Pauli group, Hilbert space, or
Hamiltonian.

A \emph{scheme Haah datum} interprets the same package on \(X=\Spec A\) by
sheaves, supports, localization, residue fibres, and nilpotent thickenings.
By convention \texttt{(ba)}, it still does not identify points of \(\Spec A\)
with translation sites.  A \emph{finite Weyl-Haah datum} adds finite
self-duality for Weyl-Heisenberg labels; AQM-81 gives the residue-field case,
while finite-ring versions require the Frobenius-ring or generating-character
sources in \path{references/heisenberg_weil/SOURCES.md}.

The most restrictive layer is the \emph{translation Haah datum}: choose an
abelian translation group \(\Lambda\) and
\[
  R=\mathbb F_p[\Lambda],
\]
for example \(R=\mathbb F_p[X^{\pm1},Y^{\pm1}]\) when
\(\Lambda=\mathbb Z^2\).  Haah's source-local package uses Laurent
coefficients for finite-support Pauli or qudit labels and exponents for
translation positions
(\path{references/lattice_stabilizer_codes/haah_1204_1063_source/alg-code-hamiltonian-v6.tex},
lines 169--183).

\subsection{The Exact Translation Package}

Let \(\Lambda\) be abelian.  Haah's qubit source takes
\[
  P=R^{2q},\qquad R=\mathbb F_2[\Lambda],
\]
as the Pauli module for finite-support Paulis on
\(\Lambda\times\{1,\ldots,q\}\), modulo phase factors; the prime-qudit
two-dimensional source uses \(R=\mathbb F_p[x^\pm,y^\pm]\)
(\path{references/lattice_stabilizer_codes/haah_1204_1063_source/alg-code-hamiltonian-v6.tex},
lines 416--420 and 532--544;
\path{references/lattice_stabilizer_codes/haah_1812_11193_source/cph2d-arXiv3.tex},
lines 321--365).

The antipode is \(\overline{g}=g^{-1}\), or \(x\mapsto x^{-1}\),
\(y\mapsto y^{-1}\); \(v^\dagger\) is transpose followed by this involution.
The trace is the identity-coefficient functional
\[
  \operatorname{tr}\left(\sum_{g\in\Lambda}a_g g\right)=a_1,
\]
not a field trace.  These are the sourced conventions in
\path{references/lattice_stabilizer_codes/haah_1204_1063_source/alg-code-hamiltonian-v6.tex},
lines 422--434 and 514--525, with the two-variable coefficient-of-\(1\)
bracket in
\path{references/lattice_stabilizer_codes/haah_1812_11193_source/cph2d-arXiv3.tex},
lines 354--365 and 505--511.

A generator matrix \(\sigma:R^t\to P\) selects a stabilizer submodule.  With
\[
  \epsilon=\sigma^\dagger\lambda_q:P\to R^t,
\]
generator commutation is \(\sigma^\dagger\lambda_q\sigma=0\), giving the
complex \(R^t\xrightarrow{\sigma}P\xrightarrow{\epsilon}R^t\) and
\(\operatorname{im}\sigma\subseteq\ker\epsilon\)
(\path{references/lattice_stabilizer_codes/haah_1204_1063_source/alg-code-hamiltonian-v6.tex},
lines 634--667).  When topological order is sourced, the stronger exactness is
\(\ker\epsilon=\operatorname{im}\sigma\)
(\path{references/lattice_stabilizer_codes/haah_1305_6973_source/alg-theory.tex},
lines 700--716;
\path{references/lattice_stabilizer_codes/haah_1812_11193_source/cph2d-arXiv3.tex},
lines 409--453).

\subsection{Why This Is Not Just \texorpdfstring{\(\Spec R\)}{Spec R}}

The central separation is role, not notation.  In AQM-90 the ring
\[
  A_T=k[x^{\pm1},y^{\pm1}]
\]
is a coordinate ring for the split torus \(T=\Spec A_T\).  In AQM-91 and
\texttt{CONVENTIONS.md} \texttt{(az)}, the ring
\[
  R_{\rm tr}=\mathbb F_p[X^{\pm1},Y^{\pm1}]
\]
is a translation group algebra acting on finite-support labels.  The same
abstract Laurent algebra can be written in both places, but the basis
\(\{X^aY^b\}_{(a,b)\in\mathbb Z^2}\) is group-basis data, not the set of
Zariski points of \(\Spec R_{\rm tr}\).

AQM-90 proves a concrete warning: as a scheme, the torus has infinitely many
closed points, while Laurent monomials have infinite full closed support and
only become Weyl operators after restriction to a finite support.  Haah's
finite-support coefficients are instead indexed by translation exponents.
Therefore no shard should read a point of \(\Spec\mathbb F_p[\Lambda]\) as a
Haah lattice site unless it has first named a finite-support bridge such as
the AQM-91 periodic bridge.

\subsection{Examples}

\paragraph{1. The infinite Laurent translation algebra.}
For \(R_{\rm tr}=\mathbb F_p[\mathbb Z^2]\), the translation basis is
\[
  \{X^aY^b:(a,b)\in\mathbb Z^2\}.
\]
This is the coefficient basis for Haah finite-support labels.  If the same
ring is viewed as a coordinate ring, then \(\Spec R_{\rm tr}\) is the split
torus studied in AQM-90.  The exponent pair \((a,b)\) is not a point of this
spectrum.

\paragraph{2. A finite periodic quotient.}
AQM-91 names the finite bridge.  Choose finite orders \(N_x,N_y\) and set
\[
  R_{G,p}=
  \mathbb F_p[X^{\pm1},Y^{\pm1}]/(X^{N_x}-1,Y^{N_y}-1)
  =\mathbb F_p[C_{N_x}\times C_{N_y}].
\]
As an \(\mathbb F_p\)-vector space it has basis
\[
  \{X^iY^j:0\le i<N_x,\ 0\le j<N_y\}.
\]
This is the finite periodic translation configuration module in AQM-91, with
\(X\) and \(Y\) acting as shifts.  It is not automatically the coordinate ring
of the finite set of sites.

\paragraph{3. The AQM-91 qubit toric square.}
For the qubit toric comparison, AQM-91 specializes to \(p=2\) and writes
\[
  P_G=R_G^4=(q_x,q_y,p_x,p_y).
\]
With the AQM-90 edge orientation, the source-local Haah matrix is
\[
  \sigma_\square=
  \begin{pmatrix}
  1+X^{-1} & 0\\
  1+Y^{-1} & 0\\
  0 & 1+Y\\
  0 & 1+X
  \end{pmatrix}.
\]
Then \(\epsilon_\square=\sigma_\square^\dagger\lambda_2\) and
\(\epsilon_\square\sigma_\square=0\) give the Haah form of the AQM-90 square
boundary identity.  AQM-91 cites the infinite-ring toric exactness
\(\ker\epsilon_{\rm tr}=\operatorname{im}\sigma_{\rm tr}\) to
\path{references/lattice_stabilizer_codes/haah_1305_6973_source/alg-theory.tex},
lines 1515--1567, and it deliberately does not assert the same finite-periodic
equality without a local finite-periodic check.

\paragraph{4. The modular thickening \( \mathbb F_3[C_3]\).}
This is a local derivation, with the nilpotent-sensitive comparison anchored
by AQM-37.  Let \(g\) generate \(C_3\).  In characteristic \(3\),
\[
  \mathbb F_3[C_3]\cong\mathbb F_3[g]/(g^3-1)
  =\mathbb F_3[g]/((g-1)^3).
\]
With \(u=g-1\), this is \(\mathbb F_3[u]/(u^3)\).  AQM-37 records the same
ring as \(\mathbb F_3[t]/(t^3-1)\cong\mathbb F_3[u]/(u^3)\), with one reduced
closed point and a thickened Hilbert layer
\(\ell^2(\mathbb F_3[u]/(u^3))\).  The translation basis is still
\(\{1,g,g^2\}\).  Thus the spectrum is a single third-order thickened point,
while the translation module has three group-basis elements.

\paragraph{5. A split product warning.}
This is a local derivation using the product-field spectrum convention of
AQM-40 and \texttt{CONVENTIONS.md} \texttt{(ad)}.  Over \(\mathbb F_3\),
\[
  \mathbb F_3[C_2]\cong \mathbb F_3[t]/(t^2-1)
  \cong \mathbb F_3[t]/(t-1)\times\mathbb F_3[t]/(t+1)
  \cong\mathbb F_3\times\mathbb F_3.
\]
Therefore its spectrum has two reduced product-field points.  The translation
basis is \(\{1,t\}\), while the product idempotents
\[
  e_+=2(1+t),\qquad e_-=2(1-t)
\]
are the split character or Fourier basis in this displayed example.  Passing
to these idempotents is extra decomposition data, not the same as remembering
the group-basis shifts used by a Haah finite-difference matrix.

\subsection{Checklist for Reaching the Translation Haah Layer}

A scheme proposal has reached the translation Haah layer only after naming or
sourcing: an abelian translation group \(\Lambda\) or finite quotient; the
group algebra \(R=\mathbb F_p[\Lambda]\) or periodic quotient; the group basis
and finite-support convention; the antipode, dagger, and identity-coefficient
trace; \(P=R^{2q}\) and \(\lambda_q\); a finite generator matrix \(\sigma\);
the excitation map \(\epsilon=\sigma^\dagger\lambda_q\) and checked isotropy
\(\sigma^\dagger\lambda_q\sigma=0\); a sourced exactness claim or the
recorded quotient \(\ker\epsilon/\operatorname{im}\sigma\); a finite periodic
or residue/Weyl bridge for finite Hilbert spaces; and an explicit warning for
any Fourier, character, residue, or \(\Spec\)-point decomposition added to the
translation basis.

\subsection{Right Lesson}

Bare \(\Spec A\) gives geometric supports, residue fields, localizations, and
nilpotent thickenings in the conventions used by AQM-28, AQM-36, and AQM-81.
It does not by itself give Haah lattice sites.  Haah's lesson for this
programme is more restrictive and more useful: once a chosen commutative ring
is a translation group algebra or a controlled quotient of one, modules over
that ring can carry symplectic duality, an antipode, an identity-coefficient
trace, and finite-difference stabilizer maps.  The physical code layer is the
extra structured package
\[
  (R,\bar{(\cdot)},\operatorname{tr},P,\lambda_q,\sigma,\epsilon)
\]
together with finite-support or finite-periodic interpretation.  A scheme
proposal may borrow this grammar, but it has not become Haah data until those
translation choices are explicit and sourced.
